\documentclass[12pt]{article}
\usepackage[utf8]{inputenc}
\usepackage[T1]{fontenc}
\usepackage[spanish]{babel}
\usepackage{amsmath, amssymb, amsfonts}
\usepackage{fancyhdr}
\usepackage{xcolor}
\usepackage[margin=3cm]{geometry}
\usepackage{graphicx}
\usepackage{float}
\usepackage{multicol}
\usepackage{enumitem}
\usepackage{tcolorbox}
\usepackage{listings}

% Configuración de listings para Python
\definecolor{codegreen}{rgb}{0,0.6,0}
\definecolor{codegray}{rgb}{0.5,0.5,0.5}
\definecolor{codepurple}{rgb}{0.58,0,0.82}
\definecolor{backcolour}{rgb}{0.95,0.95,0.92}

\lstdefinestyle{mystyle}{
    backgroundcolor=\color{backcolour},   
    commentstyle=\color{codegreen},
    keywordstyle=\color{magenta},
    numberstyle=\tiny\color{codegray},
    stringstyle=\color{codepurple},
    basicstyle=\ttfamily\small,
    breakatwhitespace=false,         
    breaklines=true,                 
    captionpos=b,                    
    keepspaces=true,                 
    numbers=left,                    
    numbersep=5pt,                  
    showspaces=false,                
    showstringspaces=false,
    showtabs=false,                  
    tabsize=2,
    language=Python
}

\lstset{style=mystyle}

% Encabezado
\pagestyle{fancy}
\fancyhf{}
\fancyhead[L]{Escuela de Ingeniería \\ Universidad de O'Higgins}
\fancyhead[R]{Programación\\ 2026-01}
\fancyfoot[C]{\thepage}

% Título comprimido
\makeatletter
\renewcommand{\maketitle}{%
  \begin{center}
    {\LARGE \@title\par}
    \vskip 0.3em
    {\large \@author\par}
    \vskip 0.3em
  \end{center}
  \vskip -0.2em  % Espacio después del título
}
\makeatother

\title{Ayudantía 5}
\author{Felipe Llach R.}

\begin{document}
\maketitle

\section{Ejercicios}

\begin{enumerate}

    \item \textbf{Cálculo de Potencia:}
    Escriba una función llamada \texttt{potencia(base, exponente=2)} que reciba una base y un exponente. Si el exponente no es proporcionado, la función debe elevar la base al cuadrado. La función debe retornar el resultado.

    \item \textbf{Estadísticas Básicas:}
    Cree una función llamada \texttt{analizar\_numeros(a, b, c)} que reciba tres números enteros. La función debe retornar tres valores: la suma de los tres, el número menor y el número mayor.

    \item \textbf{Documentación y Validación:}
    Escriba una función llamada \texttt{es\_bisiesto(year)} que determine si un año es bisiesto. 
    \begin{itemize}
        \item Debe incluir un \textbf{docstring} detallado que explique la lógica (Un año es bisiesto si es divisible por 4, pero no por 100, a menos que sea divisible por 400).
        \item Pruebe la documentación usando \texttt{help(es\_bisiesto)}.
        \item \textit{Concepto:} Docstrings y ayuda integrada.
    \end{itemize}

    \item \textbf{Simulador de Dados (Módulo \texttt{random}):}
    Importe el módulo \texttt{random} y cree una función \texttt{lanzar\_dados(cantidad)} que simule el lanzamiento de $N$ dados de 6 caras. La función debe retornar una lista con los resultados de cada dado.

    \item \textbf{Geometría Analítica (Módulo \texttt{math}):}
    Cree un módulo llamado \texttt{geometria.py} (conceptualmente o en un archivo aparte) que contenga una función \texttt{distancia\_entre\_puntos(x1, y1, x2, y2)}. Utilice \texttt{math.sqrt} y \texttt{math.pow} para calcular la distancia euclidiana entre dos puntos $(x_1, y_1)$ y $(x_2, y_2)$.

    \item \textbf{Conversor de Moneda con IVA:}
    Diseñe una función \texttt{calcular\_total(monto, tasa\_iva=0.19)} que calcule el total a pagar incluyendo el IVA. Si no se especifica la tasa, debe usar el 19\% por defecto. La función debe documentar claramente sus argumentos.

\end{enumerate}

\section*{Desafío Extra}
Cree un script principal que importe las funciones de los ejercicios anteriores y permita al usuario interactuar con ellas a través de un menú simple.

\end{document}
